A estimativa vem de \textbf{semelhan\c{c}a geométrica}: se o comprimento característico de Gulliver é $12$ vezes o de um liliputiano e ambos t\^em a mesma densidade corporal média $\rho_0$, então:
\[
\frac{M_{\rm G}}{M_{\rm L}} = \frac{\rho_0\,V_{\rm G}}{\rho_0\,V_{\rm L}} = \left(\frac{L_{\rm G}}{L_{\rm L}}\right)^3 = 12^3 = 1728.
\]

O raciocínio dos liliputianos baseia-se na hipótese de que a \textbf{necessidade alimentar diária é proporcional à massa corporal} (metabolismo proporcional à massa). Com essa hipótese:
\[
\frac{\text{alimento}_{\rm G}}{\text{alimento}_{\rm L}} = \frac{M_{\rm G}}{M_{\rm L}} = 1728.
\]

Portanto, a razão $12^3 = 1728$ é exatamente o fator de escala volumeétrica (e, portanto, de massa) entre Gulliver e um liliputiano, justificando a conclusão sobre a demanda alimentar.

